% ==================================================================
% id:       cond-013
% topic:    הסתברות מותנית ואי-תלות
% subtopic: נוסחת ההסתברות השלמה ונוסחת בייס
% solution: full
% ==================================================================
במבחן אמריקאי יש לכל שאלה 5 תשובות אפשריות. אם הסטודנטית יודעת את התשובה, היא בוחרת בה. אם לא, היא מנחשת ובוחרת אחת מ-5 התשובות באקראי. ההסתברות שהיא יודעת את התשובה לשאלה מסוימת היא 0.7. הסטודנטית ענתה נכון על השאלה. מה ההסתברות שהיא ידעה את התשובה?
\begin{solbox}
נסמן $K$ = „יודעת”, $R$ = „ענתה נכון”. נתון $\PP(K)=0.7$, $\PP(R\mid K)=1$, $\PP(R\mid K^c)=\frac15$. לפי נוסחת ההסתברות השלמה
\[
\PP(R)=1\cdot0.7+\tfrac15\cdot0.3=0.76,
\]
ולפי בייס
\[
\PP(K\mid R)=\frac{0.7}{0.76}=\frac{35}{38}\approx\ans{0.921}.
\]
\end{solbox}
